% 独立章节文件：可单独编辑，统一由 main.tex 引入。
\chapter{常温自动工况与数据采集}
\section{常温 ACC：12 拟合 + 12 验证}
2TD600 参考工况以 $(\beta,\alpha)$ 表示\textbf{（内轴角度、外轴角度）}，单位度。下列为当前\Code{production_calibration.py} 的实际列表。
% 可复用表格：修改内容后重新编译 main.tex
\begin{table}[H]\centering\small
\caption{ACC 拟合姿态（12 组）}
\begin{tabular}{@{}cc@{\hspace{12mm}}cc@{\hspace{12mm}}cc@{}}\toprule
编号 & $(\beta,\alpha)$ & 编号 & $(\beta,\alpha)$ & 编号 & $(\beta,\alpha)$\\\midrule
01 & $(0,0)$ & 05 & $(0,90)$ & 09 & $(225,135)$\\
02 & $(0,180)$ & 06 & $(180,90)$ & 10 & $(315,135)$\\
03 & $(90,90)$ & 07 & $(45,45)$ & 11 & $(45,225)$\\
04 & $(270,90)$ & 08 & $(135,45)$ & 12 & $(225,315)$\\\bottomrule
\end{tabular}\end{table}

% 可复用表格：修改内容后重新编译 main.tex
\begin{table}[H]\centering\small
\caption{ACC 独立验证姿态（12 组，不参与拟合）}
\begin{tabular}{@{}cc@{\hspace{12mm}}cc@{\hspace{12mm}}cc@{}}\toprule
编号 & $(\beta,\alpha)$ & 编号 & $(\beta,\alpha)$ & 编号 & $(\beta,\alpha)$\\\midrule
01 & $(25,35)$ & 05 & $(235,215)$ & 09 & $(220,55)$\\
02 & $(75,65)$ & 06 & $(285,245)$ & 10 & $(310,125)$\\
03 & $(125,115)$ & 07 & $(40,295)$ & 11 & $(70,235)$\\
04 & $(175,145)$ & 08 & $(140,325)$ & 12 & $(250,305)$\\\bottomrule
\end{tabular}\end{table}

六个基本方向为矩阵可观测性提供基础，其余斜姿态进一步帮助识别耦合项。每个姿态的参考重力方向应由\textbf{已确认的转台运动学与真实角度反馈}求出，而不是按按钮目标值假定完全到位。

\section{常温 GYRO：19 拟合 + 13 验证}
% 可复用表格：修改内容后重新编译 main.tex
\begin{table}[H]\centering
\caption{GYRO 速率点选择（当前试运行工艺）}
\noindent\begin{tabularx}{\linewidth}{@{}p{25mm}p{45mm}p{28mm}Y@{}}\toprule
阶段 & 参考角速度 & 工况数 & 目的\\\midrule
拟合静止 & $0^\circ/\mathrm{s}$ & 1 & 三轴零偏基准\\
拟合动态 & $\pm1,\pm2,\pm5^\circ/\mathrm{s}$，参考 X/Y/Z 各执行 & 18 & 三轴比例与耦合矩阵\\
验证静止 & $0^\circ/\mathrm{s}$ & 1 & 未参与拟合的静止残差\\
验证动态 & $\pm1.5,\pm3.5^\circ/\mathrm{s}$，参考 X/Y/Z 各执行 & 12 & 独立的速度幅值和方向检查\\\midrule
总计 &  & \textbf{32} & \textbf{19 拟合 + 13 验证}\\\bottomrule
\end{tabularx}\end{table}

实际参考轴映射由源码给出：Z 使用内轴，X 使用外轴（内轴 $0^\circ$），Y 使用外轴（内轴 $270^\circ$）；具体正方向必须通过真实转台核对。左右轴/正负转向的参考角速度应采用\textbf{实际稳定段反馈}而非简单取下发的目标值。

\section{为何低速点有用、也有边界}
若某一轴 $y=k\omega+b+\epsilon$，等噪声、带截距最小二乘下
\begin{equation}
\operatorname{Var}(\hat k)=\frac{\sigma^2}{\sum_i(\omega_i-\bar\omega)^2}.
\end{equation}
在风机塔筒等低频缓慢姿态变化应用中，$\pm1\sim\pm5^\circ/\mathrm{s}$ 是有意义的低速校准区间；但它\textbf{不能凭自身证明}中高速非线性或整机动态精度满足要求。是否增加独立高速验证，应由实际角速度需求、转台计量能力、机械行程及线缆风险共同决定。配置中的更高速建议点处于\textbf{禁用状态}，不是现有生产工况。

\section{单工况的真实执行步骤}
\begin{enumerate}
\item 检查转台通信、使能、伺服正常、反馈及时与样品配置一致。
\item 下发目标姿态并等待位置/停车反馈；角度跨 $0/360^\circ$ 时按圆周等价处理。
\item 静置等待（当前配置 3 s）；静态姿态采集 5 s；动态速率稳定后采集 5 s。
\item GYRO 拟合零速当前配置 120 s（按最新配置文件），验证零速不少于 10 s。
\item 从实时 IMU 帧取得 ACC、GYRO、芯片温度、转台姿态及反馈速率；记录工况编号与时间。
\item 对帧内及跨帧异常做稳健筛除，检查采样长度、噪声、温漂和转台稳定性；失败有上限重试。
\item 动态工况结束后下发停车，并等待停车反馈确认，才进入下一工况。
\end{enumerate}

\section{当前工程质量门限（\textbf{试运行}）}
以下值来自最新\Code{production_quality_profile.json}，与未来正式生产验收值不同；工艺文件中的\Code{approved_for_production} 目前为 \Code{false}。
% 可复用表格：修改内容后重新编译 main.tex
\begin{table}[H]\centering\small
\caption{采样质量门限（软件当前值；并非计量标准通用规定）}
\noindent\begin{tabularx}{\linewidth}{@{}Yp{32mm}Yp{32mm}@{}}\toprule
项目 & 当前值 & 项目 & 当前值\\\midrule
每工况最少有效帧 & 50 & 异常剔除比例上限 & 10\%\\
ACC 稳健剔除 & $3\times\mathrm{MAD\mbox{-}\sigma}$ & GYRO 稳健剔除 & $6\times\mathrm{MAD\mbox{-}\sigma}$\\
静止 ACC 单轴标准差上限 & 0.025 g & 静止 GYRO 单轴标准差上限 & 0.45 $^\circ$/s\\
动态 GYRO 单轴标准差上限 & 0.55 $^\circ$/s & 转台参考速率标准差上限 & 0.18 $^\circ$/s\\
转台速率均值偏离目标上限 & 0.30 $^\circ$/s & 静止反馈速率上限 & 0.05 $^\circ$/s\\
转台反馈最大年龄 & 250 ms & 常温单工况芯片温度跨度 & 1.0 $^\circ$C\\\bottomrule\end{tabularx}\end{table}

常温运行配置温度范围为 $15\sim35\, ^\circ\mathrm C$；静态采样最多允许 5 次重试、动态采样最多 3 次重试（具体计数含义以调用处为准）。质量门限用于拒绝不可用数据，\textbf{不会自动保证样本参考准确度}。
